\chapter{Giới thiệu}

\section{Bối cảnh và Động lực}

Blockchain Ethereum là một nền tảng hợp đồng thông minh phi tập trung, nơi mỗi giao dịch đều phải trả một khoản phí gọi là \textit{gas fee}. Phí này dao động liên tục theo thời gian thực, phụ thuộc vào mức độ tắc nghẽn mạng (network congestion). Trong những giờ cao điểm, chi phí gas có thể tăng gấp nhiều lần so với mức trung bình, gây ra tổn thất đáng kể cho các tổ chức và doanh nghiệp cần xử lý khối lượng giao dịch lớn định kỳ (ví dụ: các hệ thống thanh toán, sàn giao dịch phi tập trung, hoặc các ứng dụng DeFi).

Bài toán đặt ra là: \textit{Tại mỗi thời điểm, nên quyết định gửi bao nhiêu giao dịch} (trong tổng số các giao dịch đang chờ xử lý) để đồng thời: (1) tối thiểu hóa tổng chi phí gas, và (2) không vi phạm ràng buộc về thời hạn (deadline) — tức là đảm bảo tất cả giao dịch được hoàn tất trước khi hết hạn.

\section{Phạm vi Đề tài}

Đề tài này xây dựng một \textbf{hệ thống quyết định tự động} sử dụng \textbf{Offline Reinforcement Learning (Offline RL)} để giải quyết bài toán trên. Điểm đặc thù của Offline RL là agent học hoàn toàn từ \textit{dữ liệu lịch sử} (không tương tác trực tiếp với môi trường thật), phù hợp với ngữ cảnh blockchain nơi việc thử và sai trong môi trường thực là cực kỳ tốn kém.

\section{Đóng góp Chính}

\begin{itemize}
    \item Xây dựng pipeline xử lý dữ liệu thô Ethereum thành tập \textit{transition dataset} phục vụ Offline RL.
    \item Thiết kế \textbf{Hindsight Oracle} — bộ giải Quy hoạch Động (Dynamic Programming) biết trước giá gas trong tương lai — để tạo ra các quỹ đạo tối ưu làm nhãn huấn luyện.
    \item Chuyển đổi không gian hành động từ liên tục sang \textbf{rời rạc (5 bins)} nhằm giải quyết vấn đề phân phối hành động sụp đổ (policy collapse) trong IQL liên tục.
    \item Huấn luyện và đánh giá agent \textbf{Discrete CQL (Conservative Q-Learning)} trên tập dữ liệu thực tế từ Ethereum Mainnet (2024--2026).
\end{itemize}
